\textbf{结论名称：}双曲线第一定义

\textbf{适用条件：}平面上动点到两定点距离之差的绝对值为常数，且常数严格小于两定点间距（\(2a<|F_1F_2|\)）。

\textbf{变量说明：} \(F_1,F_2\) 为双曲线的两个焦点，焦距 \(|F_1F_2|=2c\)，\(a\) 为实半轴长（\(a>0\)），\(P\) 为轨迹上任意一点。

\textbf{核心公式：}
\[
\boxed{\bigl||PF_1|-|PF_2|\bigr|=2a\qquad (0<2a<|F_1F_2|)}
\]

\textbf{使用场景：} 根据定义直接判断动点轨迹是否为双曲线；在已知焦点和实轴长的条件下建立坐标系，求双曲线的标准方程；与椭圆定义对比，处理轨迹判断题。

\textbf{简单例子：} 设两定点 \(F_1(-5,0),F_2(5,0)\)，动点 \(P\) 满足 \(\bigl||PF_1|-|PF_2|\bigr|=8\)。此时 \(2a=8\Rightarrow a=4\)，\(c=5\)，且 \(2a<2c\) 成立，故轨迹为双曲线。由 \(b^2=c^2-a^2=9\) 得标准方程为 \(\frac{x^2}{16}-\frac{y^2}{9}=1\)。